%==============================================================================
% Chapitre 21 - Plateformes de tir
%==============================================================================

\section{Introduction}

Une arme n'est généralement pas utilisée seule.

Elle est intégrée à une plateforme qui assure :

\begin{itemize}
    \item son support ;
    \item son orientation ;
    \item sa stabilisation ;
    \item son transport ;
    \item son alimentation éventuelle.
\end{itemize}

La plateforme influence directement les performances globales du système
d'arme.

\begin{aretenir}

Les performances observées résultent souvent davantage de l'ensemble
arme-plateforme que de l'arme seule.

\end{aretenir}

\section{Définition d'une plateforme de tir}

Une plateforme de tir est tout système permettant de mettre en œuvre une
arme.

Elle peut être :

\begin{itemize}
    \item fixe ;
    \item mobile ;
    \item terrestre ;
    \item navale ;
    \item aérienne.
\end{itemize}

La plateforme constitue le lien entre le tireur, l'arme et l'environnement.

\section{Plateformes fixes}

Les plateformes fixes représentent le cas le plus simple.

Elles comprennent notamment :

\begin{itemize}
    \item les bancs d'essais ;
    \item les affûts ;
    \item les installations permanentes ;
    \item certaines positions fortifiées.
\end{itemize}

Elles offrent généralement :

\begin{itemize}
    \item une excellente stabilité ;
    \item une faible dispersion supplémentaire ;
    \item une bonne répétabilité des mesures.
\end{itemize}

\section{Armes portées par un tireur}

Dans ce cas, le tireur lui-même devient une partie de la plateforme.

Le système complet comprend :

\begin{itemize}
    \item le tireur ;
    \item l'arme ;
    \item les équipements associés.
\end{itemize}

Les mouvements du corps influencent directement :

\begin{itemize}
    \item le pointage ;
    \item la stabilité ;
    \item la précision.
\end{itemize}

\section{Bipieds et trépieds}

Les supports mécaniques permettent d'améliorer la stabilité.

Le bipied offre :

\begin{itemize}
    \item simplicité ;
    \item légèreté ;
    \item rapidité de mise en œuvre.
\end{itemize}

Le trépied permet généralement :

\begin{itemize}
    \item une meilleure stabilité ;
    \item un pointage plus précis ;
    \item une meilleure répétabilité.
\end{itemize}

\begin{exemple}

Une MAG58 montée sur trépied présente généralement une dispersion plus faible
qu'une MAG58 utilisée à l'épaule.

\end{exemple}

\section{Plateformes terrestres}

Les véhicules terrestres constituent aujourd'hui des plateformes majeures.

On distingue notamment :

\begin{itemize}
    \item véhicules légers ;
    \item véhicules blindés ;
    \item chars ;
    \item systèmes d'artillerie.
\end{itemize}

Ils permettent :

\begin{itemize}
    \item une mobilité importante ;
    \item une capacité d'emport élevée ;
    \item l'intégration de systèmes complexes.
\end{itemize}

\section{Influence du mouvement}

Une plateforme mobile introduit des contraintes supplémentaires.

Le véhicule peut présenter :

\begin{itemize}
    \item des translations ;
    \item des rotations ;
    \item des vibrations ;
    \item des accélérations.
\end{itemize}

Ces mouvements influencent directement :

\begin{itemize}
    \item la visée ;
    \item la précision ;
    \item les calculs balistiques.
\end{itemize}

\begin{aretenir}

Le mouvement de la plateforme modifie les conditions initiales du tir.

\end{aretenir}

\section{Tirs depuis un véhicule en mouvement}

Lorsqu'un tir est effectué depuis une plateforme mobile :

\begin{itemize}
    \item le projectile hérite initialement de la vitesse du véhicule ;
    \item la ligne de visée peut évoluer pendant le tir ;
    \item des corrections peuvent être nécessaires.
\end{itemize}

Cette situation a été abordée dans les chapitres consacrés aux référentiels.

\begin{exemple}

Deux véhicules roulant côte à côte à la même vitesse présentent un cas
particulier où certains effets relatifs disparaissent presque totalement.

\end{exemple}

\section{Tourelles}

Les tourelles permettent d'orienter l'arme indépendamment du véhicule.

Elles offrent généralement deux degrés de liberté :

\begin{itemize}
    \item azimut ;
    \item site.
\end{itemize}

Certaines architectures plus complexes ajoutent d'autres possibilités de
pointage.

\section{Stabilisation}

La stabilisation vise à maintenir l'orientation de l'arme malgré les
mouvements de la plateforme.

Les systèmes modernes utilisent :

\begin{itemize}
    \item gyroscopes ;
    \item centrales inertielles ;
    \item servomécanismes ;
    \item calculateurs numériques.
\end{itemize}

\begin{exemple}

Un char moderne peut engager une cible tout en se déplaçant grâce à la
stabilisation de son armement principal.

\end{exemple}

\section{Plateformes aériennes}

Les hélicoptères et aéronefs introduisent des contraintes supplémentaires :

\begin{itemize}
    \item vitesses élevées ;
    \item vibrations ;
    \item mouvements rapides ;
    \item variations d'altitude.
\end{itemize}

Les systèmes de conduite de tir jouent alors un rôle essentiel.

\section{Plateformes navales}

Les plateformes navales doivent également compenser :

\begin{itemize}
    \item roulis ;
    \item tangage ;
    \item lacet ;
    \item mouvements dus à la mer.
\end{itemize}

Ces mouvements peuvent être importants même lorsque la cible est immobile.

\section{Plateformes et balistique}

La plateforme influence plusieurs paramètres balistiques :

\begin{itemize}
    \item vitesse initiale relative ;
    \item orientation initiale ;
    \item dispersion ;
    \item stabilité du pointage ;
    \item précision globale.
\end{itemize}

Dans certains cas, l'influence de la plateforme peut dépasser celle de
l'arme elle-même.

\section{Plateformes et simulation}

Dans les simulations modernes, la plateforme est souvent modélisée comme un
système indépendant relié :

\begin{itemize}
    \item au moteur balistique ;
    \item au système de visée ;
    \item au moteur physique ;
    \item au modèle de mobilité.
\end{itemize}

Cette approche permet de représenter :

\begin{itemize}
    \item les vibrations ;
    \item les mouvements ;
    \item les stabilisations ;
    \item les limitations mécaniques.
\end{itemize}

\begin{simulation}

Dans les simulateurs militaires, la fidélité du modèle de plateforme est
souvent aussi importante que celle du modèle balistique.

\end{simulation}

\section{Vers les systèmes de conduite de tir}

Les plateformes modernes sont généralement associées à des systèmes
électroniques chargés d'assister le tireur.

Ces systèmes réalisent notamment :

\begin{itemize}
    \item des mesures ;
    \item des calculs balistiques ;
    \item des corrections ;
    \item des fonctions de stabilisation.
\end{itemize}

Ils seront étudiés dans le chapitre suivant.

\section{À retenir}

\begin{aretenir}

\begin{itemize}
    \item La plateforme constitue le support de l'arme.
    \item Les mouvements de la plateforme influencent les conditions du tir.
    \item Les tourelles permettent le pointage de l'armement.
    \item Les systèmes de stabilisation compensent les mouvements parasites.
    \item Les plateformes modernes sont étroitement liées aux systèmes de
    conduite de tir.
\end{itemize}

\end{aretenir}
